\subsection{Superposition}

We have already seen that:
\begin{enumerate}
    \item \textbf{A qubit is represented as a vector}:
    \begin{equation*}
        \ket{\psi} = a\ket{0} + b\ket{1}
    \end{equation*}
    Where $a$ and $b$ are complex numbers, and $\ket{0}$ and $\ket{1}$ are the basis states of the qubit.


    \item \textbf{Measurement probabilities depend on the chosen basis}:
    \begin{itemize}
        \item In the computational basis $\{\ket{0}, \ket{1}\}$, the probabilities of measuring $\ket{0}$ and $\ket{1}$ are $|a|^2$ and $|b|^2$, respectively.
        \item In another basis, such as the Hadamard basis $\{\ket{+}, \ket{-}\}$, the probabilities of measuring $\ket{+}$ and $\ket{-}$ depend on different coefficients, which can be calculated by expressing $\ket{\psi}$ in terms of the new basis states.
    \end{itemize}

    
    \item \textbf{The same state can lead to deterministic or probabilistic outcomes depending on the measurement basis}. For example, if $\ket{\psi} = \ket{0}$, measuring in the computational basis always yield $\ket{0}$ (deterministic):
    \begin{equation*}
        \ket{0} = 1\ket{0} + 0\ket{1} = \ket{0}
    \end{equation*}
    Whereas measuring in the Hadamard basis yields $\ket{+}$ or $\ket{-}$ with equal probability (probabilistic):
    \begin{equation*}
        \ket{0} = \frac{1}{\sqrt{2}}(\ket{+} + \ket{-})
    \end{equation*}
\end{enumerate}
These observations raise an essential question: ``\emph{why can a single state exhibit such different behaviors under different measurements?}''. The answer is \textbf{superposition}.

\highspace
\begin{flushleft}
    \textcolor{Green3}{\faIcon{question-circle} \textbf{What is Superposition?}}
\end{flushleft}
\definition{Superposition} is a fundamental principle of quantum mechanics that \textbf{applies to all quantum systems}, not just qubits.
\begin{definitionbox}[: Superposition]
    \textbf{Superposition} means that a \textbf{system can exist in several possible states at the same time until a measurement is made}. It is \textbf{basis-dependent}, meaning a state may be a superposition in one basis but not in another, for example:
    \begin{equation*}
        \ket{v} = \frac{1}{\sqrt{2}}\ket{0} + \frac{1}{\sqrt{2}}\ket{1}
    \end{equation*}
    Is a superposition with respect to the computational basis $\{\ket{0}, \ket{1}\}$, but it is \textbf{not} a superposition in the Hadamard basis $\{\ket{+}, \ket{-}\}$, where it can be written as:
    \begin{equation*}
        \ket{v} = 1\ket{+} + 0\ket{-} = \ket{+}
    \end{equation*}

    \highspace
    Formally, let $\left\{\ket{u}, \ket{u^{\perp}}\right\}$ be an orthonormal basis of a single-qubit Hilbert space. A qubit is said to be in \textbf{superposition} with respect to this basis if it can be written as:
    \begin{equation*}
        \ket{\psi} = a\ket{u} + b\ket{u^{\perp}}
    \end{equation*}
    Where both coefficients $a$ and $b$ are non-zero and satisfy $\left|a\right|^{2} + \left|b\right|^{2} = 1$ (normalization condition). If one of the coefficients is zero, the state coincides with one of the basis vectors and the measurements in that basis becomes deterministic, thus it is not a superposition.
\end{definitionbox}

\highspace
More in general, a quantum state $\ket{\psi}$ in a system with multiple possible states can exist in a \textbf{linear combination} of these states:
\begin{equation*}
    \ket{\psi} = c_{1} \ket{\psi_{1}} + c_{2} \ket{\psi_{2}} + \dots + c_{n} \ket{\psi_{n}}    
\end{equation*}
Where:
\begin{itemize}
    \item $\ket{\psi_{1}}, \ket{\psi_{2}}, \dots, \ket{\psi_{n}}$ are \textbf{basis states} of the system.
    
    \item $c_{1}, c_{2}, \dots, c_{n}$ are \textbf{complex probability amplitudes}.

    \item The system is in \textbf{all states at the same time}, with the probability of measuring each state given by $\left|c_{i}\right|^{2}$.

    \item The state is \textbf{normalized}, meaning:
    \begin{equation*}
        \left|c_{1}\right|^{2} + \left|c_{2}\right|^{2} + \dots + \left|c_{n}\right|^{2} = 1
    \end{equation*}
\end{itemize}

\begin{examplebox}[: Single-Particle Systems (Quantum Mechanics)]
    In general quantum mechanics, a particle can exist in \textbf{multiple positions simultaneously} as a wave function $\Psi(x)$:
    \begin{equation*}
        \ket{\Psi} = \displaystyle\int \Psi(x) \: \ket{x} \: \mathrm{d}x  
    \end{equation*}
    \begin{itemize}
      \item The particle is in a \textbf{superposition of all possible positions} $\ket{x}$.
      \item Measurement collapses the wave function to a single position.
    \end{itemize}
\end{examplebox}

\noindent
\textcolor{Green3}{\faIcon{question-circle} \textbf{How Can a Particle Exist in Multiple States at Once?}} This question touches the heart of quantum mechanics, where our everyday intuition breaks down. The \textbf{key idea} is that a \textbf{quantum particle} is not just a tiny ball, it \textbf{is a wave function that spreads across multiple possibilities at once}.
\begin{enumerate}
    \item \important{Quantum Particles Are Waves, Not Just Points}. In classical physics, we think of particles as tiny, solid objects, like a small ball that always has a precise position and velocity.
    
    In quantum mechanics, however, particles behave more like waves. These waves are described by a wave function $\Psi(x)$, which represents the \textbf{probability of finding the particle at different locations}.
    
    The key idea is:
    \begin{itemize}
        \item The \textbf{wave function spreads across space}, meaning the \textbf{particle does not have a single location} before measurement.
        \item Instead, it exists in a superposition of all possible locations.
    \end{itemize}


    \item \important{The Double-Slit Experiment: Proof That a Particle Can Be in Two Places at Once}. The double-split experiment was first performed by Thomas Young in 1801 to demonstrate the \textbf{wave nature of light}. Later, in the 20th century, it was repeated with \textbf{electrons}, \textbf{atoms}, and even \textbf{large molecules}, revealing that \textbf{matter itself exhibits wave-like behavior}. The classical expectation for the double-slit experiment is:
    \begin{itemize}
        \item \textbf{Classical Particles} (e.g., tiny balls). Imagine firing small balls at a barrier with two slits:
        \begin{itemize}
            \item Each ball passes through \textbf{either} slit 1 \textbf{or} slit 2.
            \item On a screen behind the slits, we would observe \textbf{two bright bands}, corresponding to the two slits.
        \end{itemize}
        The result is a simple pattern of two lines, no interference.

        \item \textbf{Classical Waves} (e.g., Water or Light). If instead we send \textbf{waves} (like water waves or classical light):
        \begin{itemize}
            \item Each wave passes through \textbf{both slits simultaneously}.
            \item The waves emerging from the slits \textbf{interfere}: \textbf{constructive interference} creates bright bands, while \textbf{destructive interference} creates dark bands.
        \end{itemize}
        The result is an \textbf{interference pattern} of alternating bright and dark bands.
    \end{itemize}
    The quantum reality is even more surprising:
    \begin{itemize}
        \item \textbf{Quantum Particles} (e.g., photons or electrons). When the experiment is performed with \textbf{single quantum particles} (photons or electrons), something remarkable happens:
        \begin{enumerate}
            \item \textbf{Particles are emitted one at a time}.
            \item Each particle is detected as a \textbf{localized point} on the screen.
            \item After many particles are detected, an \textbf{interference pattern gradually emerges}.
        \end{enumerate}
        This means that each particle behaves as if it passes through \textbf{both slits simultaneously}. The particle is described by a \textbf{wave function}, representing a \textbf{superposition of paths}:
        \begin{equation*}
            \ket{\psi} = \frac{1}{\sqrt{2}}\left(\ket{\text{slit 1}} + \ket{\text{slit 2}}\right)
        \end{equation*}
        The interference pattern arises from the \textbf{coherent combination} of these two possibilities.
    \end{itemize}
    Suppose we place detectors at the slits to determine \textbf{which slit} the particle passes through. The interference pattern \textbf{disappears} and the screen now shows \textbf{two bands}, just like classical particles. This is because measuring the particle's path \textbf{entangles} (i.e., correlates) it with the environment (the detector), leading to \textbf{decoherence} (loss of quantum coherence). The superposition:
    \begin{equation*}
        \frac{1}{\sqrt{2}}\left(\ket{\text{slit 1}} + \ket{\text{slit 2}}\right)
    \end{equation*}
    Becomes effectively a \textbf{classical mixture}:
    \begin{equation*}
        \frac{1}{2}\ket{\text{slit 1}}\bra{\text{slit 1}} + \frac{1}{2}\ket{\text{slit 2}}\bra{\text{slit 2}}
    \end{equation*}
    Without coherence, \textbf{interference cannot occur}.
    
    \highspace
    In summary, the double-slit experiment demonstrates the principle of quantum superposition. When quantum particles such as photons or electrons are sent one at a time through two slits, they form an interference pattern on a detection screen. This indicates that each particle behaves as if it travels through both slits simultaneously, corresponding to the superposition:
    \begin{equation*}
        \ket{\psi} = \frac{1}{\sqrt{2}}\left(\ket{\text{slit 1}} + \ket{\text{slit 2}}\right)
    \end{equation*}
    If a measurement is performed to determine which slit the particle passes through, the interference pattern disappears due to decoherence, and the system behaves like a classical probabilistic mixture.
    \begin{center}
        \href{https://youtu.be/A9tKncAdlHQ}{\emph{Double Slit Experiment explained! by Jim Al-Khalili}}\hspace{2em}\qrcode{https://youtu.be/A9tKncAdlHQ}
    \end{center}


    \item \important{Quantum Superposition: More Than Just Probability}. A common misconception is that a particle in superposition is just \textbf{an unknown state}, like a coin that is either heads or tails, but we just don't know which. This is wrong, because quantum superposition is much deeper.

    A quantum state is a \textbf{combination of all possibilities}. \emph{Until measurement}, the system is in \textbf{all possible states at once}.
    
    Mathematically, for an electron in \textbf{two locations} $x_{1}$ and $x_{2}$:
    \begin{equation*}
        \ket{\psi} = a \ket{x_{1}} + b \ket{x_{2}}
    \end{equation*}
    \begin{itemize}
        \item The electron is \textbf{literally in both places simultaneously}.
        \item The coefficients $a$ and $b$ are \textbf{complex numbers representing the probability amplitudes}.
        \item \textbf{Interference} between these amplitudes \textbf{creates quantum effects that cannot be explained by classical probability}.
    \end{itemize}


    \item \important{Superposition in Quantum Computing}. Quantum computing \textbf{directly uses} the fact that a particle can be in multiple states at once.

    \begin{itemize}
        \item A \textbf{classical bit} can only be \textbf{0 or 1}.
        \item A \textbf{qubit (quantum bit)} can be in a superposition:
        \begin{equation*}
        \ket{\psi} = a \ket{0} + b \ket{1}
        \end{equation*}
    \end{itemize}
    This means a quantum computer can \textbf{perform many calculations simultaneously}. Therefore, \textbf{superposition allows quantum computers to process information exponentially faster than classical computers for certain tasks}.

    
    \item \important{Why Don't We See Superposition in Everyday Life?} In our daily experience, objects are \textbf{not in multiple states at once} because of a process called \textbf{quantum decoherence}.

    Quantum superposition is \textbf{fragile}. When a quantum system interacts with the environment (air, light, etc.), the superposition \textbf{collapses into one definite state}. This is why large objects (like humans or cars) do not appear in multiple places at once.
    
    However, experiments (like the Double-Slit experiment) confirm that superposition is real at microscopic scales (electrons, photons, atoms, and even molecules).
\end{enumerate}
The key \textbf{properties of Superposition} are:
\begin{enumerate}
    \item \textbf{Linearity}: Any combination of valid quantum states is also a valid quantum state.
    \item \textbf{Interference}: Quantum states in superposition can interfere, leading to constructive or destructive interference.
    \item \textbf{Measurement Collapse}: When measured, the superposition collapses into a single outcome.
    \item \textbf{Phase Information}: Unlike classical probabilities, quantum superpositions include complex phases that affect interference patterns.
\end{enumerate}